\subsection{Purpose}

This appendix fixes a draft formalization of task-driven context assembly as a mechanism for assembling context not from the entire dialogue history, but from the affected objects of the structural realization.

\subsection{Goal}

The goal of the mechanism is to obtain a compact and relevant context for the next task by using $R = \langle G, \Sigma, C \rangle$ as the single source of structural localization.

\subsection{Input}

The input consists of: a new task $q$; the current structural realization $R = \langle G, \Sigma, C \rangle$; and a history of decisions $H$, containing previous clarifications, accepted assumptions, and unresolved questions.

\subsection{Localizing the task in $R$}

First, the task $q$ is localized in the structural realization. The result of this step is the initial set of affected objects:


$$A_0(q) = F_0 \cup T_0,$$

where $F_0$ are the functions explicitly mentioned or semantically corresponding to the task, and $T_0$ are the tables explicitly mentioned or directly connected to the functions found.

In practice, this means anchoring the task to the existing graph nodes, their contracts, and their signatures, rather than to the chat history as a whole.

\subsection{Choosing the local neighborhood}

After localization, a local structural neighborhood $N(A_0)$ is built. In the minimal version, it includes:

\begin{itemize}
\item the affected functions and tables themselves;
\item the tables from $C_{\mathrm{in}}$, $C_{\mathrm{out}}$, $R$, $W$ for the functions found;
\item the functions from $T_{\mathrm{in}}$, $T_{\mathrm{out}}$;
\item neighboring nodes at graph distance one.
\end{itemize}

The idea is simple: not the entire graph enters the context, but only the nearest structural zone in which the change can realistically affect system behavior.

\subsection{Lifting history by structural keys}

After $N(A_0)$ has been built, the history $H$ is lifted not in full, but only by the structural keys from the neighborhood found. The following are selected:

\begin{itemize}
\item past decisions explicitly related to functions and tables from $N(A_0)$;
\item previously recorded blocking or advisory remarks concerning the same objects;
\item accepted assumptions that have not yet been materialized in the current version of $R$.
\end{itemize}

Thus, history ceases to be an independent source of truth and becomes an auxiliary layer indexed by the objects of the structural realization.

\subsection{Assembling task-specific context}

The resulting task-specific context $\mathrm{CTX}(q)$ consists of three parts: the current state of the affected objects; the local structural neighborhood; and the relevant historical decisions lifted by the same keys. That is,


$$\mathrm{CTX}(q) = \mathrm{Current}(A_0 \cup N(A_0)) \cup \mathrm{History}(A_0 \cup N(A_0)).$$

Unrelated graph fragments and unrelated history are not included in the context, even if they are topically similar to the task.

\subsection{Output}

The output is a compact structurally anchored context suitable for: modifying $R$; repeated verification; and preparing the next model step.

\subsection{Illustrative example}

If the task is ``add error logging for registration,'' the mechanism localizes it in the registration function, its error path, and the nearest logging sink. The history is then lifted by those same objects: it may turn out that a prior decision already existed to log errors separately from the success path, but that decision was never introduced into the current realization. The task-specific context therefore contains precisely this local zone rather than the full conversation and rather than the full system.

\subsection{Status}

In the current paper, task-driven context assembly is introduced as an architectural principle and a draft formalization. The full algorithm for neighborhood selection, ranking of historical fragments, and evaluation of context sufficiency remains future work.
